\section{Introduction}
\label{sec:intro}

A gateway that terminates HTTP traffic for web front-ends and device APIs is the natural place to inspect requests at the industrial edge. Most deployed inspection is signature-based~\cite{modsecurity,roesch1999,suricata}: rules are curated by hand and generalise poorly to polymorphic or unseen payloads. A learned classifier would adapt, but it needs traffic, and in an Industrial Data Space such as Gaia-X~\cite{gaiax2021} the requests are the very thing that cannot leave the site: they embed proprietary API queries, tokens and process telemetry. Federated learning~\cite{mcmahan2017} avoids moving the data by training on each gateway and exchanging only model parameters, so that a site can benefit from an attack first seen elsewhere without any request leaving either site.

Whether that benefit materialises at a given site is less obvious than the framing suggests. Federated averaging can converge to a worse solution than local training when client objectives diverge~\cite{li2020convergence}, and what ``diverge'' means for web traffic is specific. Two sites may face different attack campaigns against the same kind of application, or the same attacks against applications whose benign requests look nothing alike, a web front-end with path words against an API with JSON endpoints. These are different kinds of non-IID data, and a classifier that must separate benign from malicious is affected by them differently. We found no study that separates the two for a federated WAF, although the search behind that statement was bounded (Section~\ref{sec:background}).

Our question is therefore: \emph{when does federating a lightweight HTTP payload classifier across sites with different benign profiles or attack mixes help, and at what cost and security limit?} We study a deliberately small model, a logistic regression over hashed character trigrams with a fixed-size update, because it fits an ARM gateway and because its simplicity keeps the comparison between federated and isolated training interpretable. $K$ gateways each hold traffic that must stay on site, and one cloud aggregator combines their parameters. Evaluation is on unique-payload-disjoint splits, at a matched false-positive budget, over eight seeds. The prototype is a payload-scoring component for such a gateway. It does not identify or authenticate the requester, so it is one input to an access decision and not a zero-trust architecture.

\subsection{Contributions}
\begin{enumerate}
    \item \textbf{A measurement of when federation helps.} Over two constructed heterogeneity axes and two feature sets, federation lifts the weaker silo and reduces cross-domain false positives, matches the best silo on aggregate ROC-AUC within seed variation when sites differ in benign profile, and does not reliably exceed the best silo on SQL-injection recall (Section~\ref{sec:evaluation}). We also report what an operating-point-aware evaluation of the gateway on CSIC~2010 shows: the deployed scoring rule flags all normal requests in that corpus.
    \item \textbf{A prototype authenticated robust aggregator.} Updates carry HMAC-SHA256 tags with timestamps and single-use nonces, are checked against an enrolment list, and are aggregated by coordinate-wise median after dimension and norm checks. In simulation with 2 of 7 clients sign-flipping, the median retains ROC-AUC 0.984 against 0.263 for averaging. The threat model, the controls and their limits, including a plain-HTTP live channel with ML-KEM benchmarked only as a component, are set out in Sections~\ref{sec:system} and~\ref{sec:security}.
    \item \textbf{Cost and dataset-construction lessons.} A fixed 1{,}040-byte update reduces cloud-bound traffic by 98.88\% at the deployed 15\,s interval (75.0\% over a 20-round training run, where every round and both directions are charged). Two construction errors, resampling before splitting on unique payloads and rendering the two classes with different request templates, inflate results on this task and are documented with their effect.
\end{enumerate}
What remains open is how far a result on two clients and public payload lists can be pushed, and which parts of earlier work already answer pieces of the question.
